\documentclass[10pt,a4paper]{amsart}
\usepackage[margin=19mm]{geometry}
\usepackage{amsmath,amssymb,mathtools}
\usepackage[scr=rsfs,cal=euler]{mathalfa}
\usepackage{xfrac}
\usepackage[unicode,hidelinks,bookmarks=false]{hyperref}
\newcommand{\ie}{i.e.\ }
\newcommand{\cf}{cf.\ }
\newcommand{\resp}{respectively\ }
\newcommand{\Subsection}{\subsection}
\providecommand{\nobreakdash}{\nobreak\hbox{-}}
\setlength{\emergencystretch}{2em}
\multlinegap0pt
\usepackage{amssymb}
\usepackage{dsfont}
\usepackage[shortlabels]{enumitem}
\newtheorem{remark}{Remark}
\newtheorem{theorem}{Theorem}
\newtheorem{proposition}{Proposition}
\newtheorem{lemma}{Lemma}
\def\A{{\bf A}}
\def\C{{\bf C}}
\def\GL{\mathbf{GL}}
\def\GSp{{\mathbf{GSp}}}
\def\GU{{\mathbf{GU}}}
\def\det{\mathrm{det}}
\title{On Periods and $L$-functions for $\GL_4 \times \GL_2$}
\author{Antonio Cauchi, Armando Gutierrez Terradillos}
\date{Source: arXiv:2602.14586v2 (CC BY 4.0)}
\begin{document}
\maketitle

\noindent\emph{Source and licence.} On Periods and $L$-functions for $\GL_4 \times \GL_2$. Version arXiv:2602.14586v2 (CC BY 4.0), distributed by arXiv under the Creative Commons Attribution 4.0 International licence, http://creativecommons.org/licenses/by/4.0/, as recorded in the arXiv licence metadata for this exact version and shown on the abstract page https://arxiv.org/abs/2602.14586v2. Full licence notice: \texttt{licenses/arxiv-2602.14586-CC-BY-4.0.txt}.\par\smallskip

\noindent\textit{Editorial scope.} Counted unit: the article's proposition Unfolding:Spherical:Period (source0.tex lines 2141--2146). The article's complete original proof of it is delivered with it (source0.tex lines 2147--2177, one \texttt{proof} environment of 31 lines). No mathematics is written, repaired or re-derived: the statement and the proof are the article's own lines, in the article's own notation, and no retained line is substituted or re-flowed.  Retained uncounted as the source-local context the proof needs: lines 482--486; lines 511--518; lines 1475--1491; lines 1963--1966; lines 2026--2038; lines 2096--2096; lines 2100--2104.  Nothing was omitted from any retained span, so the package carries no omission record.  No retained line contains a citation, so the wrapper carries no bibliography and no bibliography entry.  No retained line contains a figure environment, an image-inclusion command or drawing code, so no image is delivered and the package declares no asset dependency.  Editorial formatting only, in addition: an AMS \texttt{amsart} preamble that restates the article's own declarations and macros used by the retained lines, namely the environments \texttt{remark}, \texttt{theorem}, \texttt{proposition}, \texttt{lemma} on separate counters, together with the standard packages those lines need.  The retained environments print in this document as Remark 1, Theorem 1, Proposition 1, Lemma 1 in order of appearance, whereas the article prints the same items with the numbering of its own full document.  The complete native source of the article as distributed in the arXiv e-print for this exact version is bundled unchanged as \texttt{sources/194723/source0.tex}.
\begin{remark}\label{remark_stupid_on_Siegel_parabolic_so22}
For instance, when $r = 2$ and $V = \mathbb{H}^2$ with quadratic form associated with $I'_4$, then $Q_{(2,0)} = L_{(2,0)}\, N_{(2,0)} $ is such that 
\[L_{(2,0)} = \left \{ \left(\begin{smallmatrix} g  & \\ &\mu I_2'\;^{t}g^{-1}I_2'\end{smallmatrix}\right), \,\, g \in \GL_2, \, \mu \in \GL_1 \right \},\]
\[N_{(2,0)} = \left \{ \left(\begin{smallmatrix} 1 & & x & \\ & 1 & & -x\\ & & 1 & \\ & & & 1 \end{smallmatrix}\right), \,\, x \in \mathbb{G}_a \right \}.\]
\end{remark}

\begin{remark}\label{Remark:Siegel:Parabolic:Exceptional}
    Note that $\rho^{-1}Q_{(2,0)}$, with $Q_{(2,0)}$  the Siegel parabolic of $\mathbf{GSO}_{\mathbb{H}^2}$, corresponds to the parabolic subgroup $(B_{\GL_2}^{\rm opp} \times \GL_2) /\{ (z,z)\}$ of $(\GL_2 \times \GL_2)/\{ (z,z)\}$. Indeed, with the identifications of Remark \ref{remark_stupid_on_Siegel_parabolic_so22}, we have
\[\rho \left( \left(\begin{smallmatrix}
    1 &  \\  & {\mu}^{-1} \cdot {\rm det}(g) 
\end{smallmatrix} \right), g \right) =  \left(\begin{smallmatrix} g  & \\ &\mu I_2'\;^{t}g^{-1}I_2'\end{smallmatrix}\right),\, \,  \rho \left( \left(\begin{smallmatrix}
    1 &  \\ x & 1 
\end{smallmatrix} \right), I_2 \right) = \left(\begin{smallmatrix} 1 & & x & \\ & 1 & & -x\\ & & 1 & \\ & & & 1 \end{smallmatrix}\right). \]
\end{remark}

\begin{theorem}\label{Final:Theorem:Central:Value}
    Let $\pi$ be a globally generic cuspidal automorphic representation of either $\GL_4(\A)$ or $\GU_{2,2}(\A)$ and let $\sigma$ be a cuspidal automorphic representation of $\GL_2(\A)$ such that as a Hecke character on $F^\times \backslash \A^\times$ \begin{equation*}
    \omega_\pi \omega_\sigma = 1.
\end{equation*}
We also let $S$ be a finite set of places containing the archimedean places for $F$ and the ramified places for $\pi$, $\sigma$, and $E/F$. Let $\varphi = \otimes_v \varphi_v$ and $\Psi = \otimes_v \Psi_v$ be cusp forms in $\pi$ and $\sigma$, and let  $\phi_{s} = \otimes_v \phi_{v,s} \in I_{P_{(3,0)}}(s,\mathbf{1})$. If 
$$ L^S(1/2, \pi \otimes \sigma, \wedge^2 \otimes  \mathrm{std}_2) \prod_{v \in S} Z\left(1/2,W_{\pi_v}, W_{\sigma_v,f_{\phi_{1/2,v}}^{\Psi_v}} \right) \ne 0,$$ then either 
\begin{align}
\mathcal{P}_{\mathbf{J}}( \varphi \otimes \Psi) : = \int_{Z_{\GSp_6}(\A)  \backslash[ \mathbf{J}]} \Psi(g_1) \varphi\left( \left(\begin{smallmatrix}
    a& &b\\ &g_2& \\ c& &d
\end{smallmatrix}\right) \right )  dg_1 \, dg_2 \ne 0,  
\end{align}
 where we denoted $g_1 = \left(\begin{smallmatrix}a&b\\c&d\end{smallmatrix}\right)$, or 
\begin{align}\label{FormulaWithTheta}
    \int_{Z_{\GSp_6}(\A)\backslash [\GL_2\boxtimes \GSp_4]} \int_{[\mathbf{O}(D)]} \theta(h h_1, \iota( g_1  , g_2) ; f^D) \Psi^w( g_1 ) \varphi(g_2) \, d h \, dg_1\,dg_2 \ne 0,
\end{align}
where $\mathbf{O}(D)$ is the orthogonal group attached to the norm form of a non-split quaternion algebra $D/F$ and $f^D \in \mathcal{S}(D(\A)^3)$ is such that ${\rm pr}_{\Pi(D)} \phi_{1/2}  = \lambda_D(f^D)^{\sim}$. 
\end{theorem}

\begin{align}
    \label{Embedding:GL1}\GL_1&\to \GL_4\times \GL_1,\\
    a&\mapsto (aI_4,a^2).\nonumber
\end{align}

\begin{proposition}\label{Proposition:Pullback:Orthogonal}
    Let $\sigma \subset \mathcal{A}_{\rm cusp}([\GL_2])$, $\Psi',\Psi''$ be cusp forms in the space of $\sigma$, and let $\Phi_{s} \in  I_{Q_{(6,0)}}(s,\mathbf{1}) $ be a holomorphic section.
    \begin{enumerate}
    \item For $\mathrm{Re}(s)>>0$, we have
    \begin{align*}
    q_{\overline{\Psi}'\boxtimes\Psi''}(\Phi_{s})(x) := \int_{\mathbf{SO}_{\mathbb{H}^2}(\A)}(\overline{\Psi}'\boxtimes\Psi'')(ht_{\nu'(x)})\Phi_{s}(q_3\kappa(ht_{\nu'(x)},x))dh\in I_{Q_{(2,4)}}(s,\overline{\sigma}\boxtimes \sigma).
    \end{align*}
    \item We obtain the pullback formula
    \[\int_{[\mathbf{SO}_{\mathbb{H}^2}]}(\overline{\Psi}'\boxtimes\Psi'')(ht_{\nu'(x)})\mathrm{E}_{Q_{(6,0)}}(\kappa(h t_{\nu'(x)},x),\Phi_{s})dh = \mathrm{E}_{Q_{(2,4)}}(x,q_{\overline{\Psi}'\boxtimes\Psi''}(\Phi_{s})).\]
    \item Let $S$ be a finite set of places containing the archimedean places and the ramified places of $\sigma$ and $\Phi_s$. If $\overline{\Psi}'\boxtimes\Psi'' \in   \sigma\boxtimes \overline{\sigma}$ and $\Phi_{s} = \otimes_v \Phi_{s,v}\in I_{Q_{(6,0)}}(s,\mathbf{1})$ is normalized so that $\Phi_{s,v}(1) = 1$ when $v \not \in S$, we have \[q_{\overline{\Psi}'\boxtimes\Psi''}(\Phi_{s})(I_8) =  \frac{L^S(5/2+s,\sigma\boxtimes\overline{\sigma},\mathrm{std})}{\zeta^S(5+2s)\zeta^S(3+2s)} \int_{\mathbf{SO}_{\mathbb{H}^2}(\A_S)}(\overline{\Psi}'\boxtimes\Psi'')(h)\Phi_{s,S}(q_3\kappa(h,I_8))dh,\] 
    where $\A_S = \prod_{v\in S} F_v$ and $ \Phi_{s,S} = \otimes_{v \in S} \Phi_{s,v}$. In particular $s\mapsto q_{\overline{\Psi}'\boxtimes\Psi''}(\Phi_{s})$ admits meromorphic continuation to all $s\in \C$. Hence,  $q_{\overline{\Psi}'\boxtimes\Psi''}(\Phi_{s})\in I_{Q_{(2,4)}}(s,\overline{\sigma}\boxtimes \sigma)$ for every $s$ at which it is holomorphic.
    \end{enumerate}
\end{proposition}

\begin{equation}\label{Integral:Section:Period:Final}\int_{[ \GL_1 \backslash (\GL_4 \times \GL_1)]}\varphi(r)\mathrm{E}_{Q_{(2,4)}}(e(r,a),q_{\overline{\Psi}'\boxtimes\Psi''}(\Phi_{s})) drda,\end{equation}

\begin{lemma}\label{Final:Double:Quotient}
    Let $Q_{(2,4)}$ be the parabolic subgroup of $\mathbf{GSO}_{\mathbb{H}^4}$ with Levi component equal to $\GL_2\times \mathbf{GSO}_{\mathbb{H}^2}$. We have that 
    \[Q_{(2,4)}(F)\setminus \mathbf{GSO}_{\mathbb{H}^4}(F)/e(\GL_4\times \GL_1)(F) = \{\omega_i\}_{i = 1,\cdots,7}.\]
    For $1\leq i \leq 6$, the stabilizer of $\omega_i$ in $\GL_4$ contains a normal unipotent subgroup. Moreover, the stabilizer of the (open) orbit $ \omega_7$ is isomorphic to   \[\left\{\left(\left(\begin{smallmatrix}g_1& \\ &g_2\end{smallmatrix}\right),\mathrm{det}(g_1)\right),\;g_1,g_2\in \GL_2(F)\right\}\subset \GL_4(F)\times \GL_1(F).\]
\end{lemma}

\begin{proposition}\label{Unfolding:Spherical:Period} When $\mathrm{Re}(s)>>0$, the integral \eqref{Integral:Section:Period:Final} equals 
      \begin{align*}
        \int_{ \mathbf{Y}(\A)^\circ } \Lambda_{r,a,s}'(I_2)\int_{Z_{\GL_4}(\A)\backslash [\GL_2\times \GL_2]} \!\!\!\!\varphi({\rm diag}(g_1,g_2) r) \Lambda''_{r,a,s}(g_2) dg_1dg_2drda.  
             \end{align*}
   where  $\mathbf{Y}(\A)^\circ :=  (\GL_2 \times \GL_2)(\A) \backslash (\GL_4 \times \GL_1)(\A)$.
\end{proposition}

\begin{proof}
    We proceed as in Theorem \ref{Final:Theorem:Central:Value} by unfolding the integral \eqref{Integral:Section:Period:Final}. For ${\rm Re}(s)$ big enough, the Eisenstein series unfolds and, by Lemma \ref{Final:Double:Quotient} and the cuspidality of $\varphi$ and $\Psi'\boxtimes\Psi''$, \eqref{Integral:Section:Period:Final} is shown to be equal to    \begin{equation}\label{First:Eq:Prop:Final:Unfolding}\int_{ (\GL_2 \times \GL_2)(F) \GL_1(\A) \backslash (\GL_4 \times \GL_1)(\A)}\varphi(r) q_{\overline{\Psi}'\boxtimes\Psi''}(\Phi_{s})(\omega_7 e(r,a)) drda,\end{equation}  
    where we recall that $\GL_1$ embeds into $\GL_4\times \GL_1$ via \eqref{Embedding:GL1}. Collapsing the integral over $[\GL_2 \times \GL_2]$, we get 
\begin{equation}\label{Second:Eq:Prop:Final:Unfolding}\int_{ \mathbf{Y}(\A)^\circ } \int_{Z_{\GL_4}(\A)\backslash [\GL_2\times \GL_2]} \!\!\!\!\!\!\!\!\!\!\! \varphi({\rm diag}(g_1,g_2)r) q_{\overline{\Psi}'\boxtimes\Psi''}(\Phi_{s})(\omega_7 e({\rm diag}(g_1,g_2) r,\det(g_1) a)) drda.\end{equation} 
    Consider the element
    \[\alpha := \left(\begin{smallmatrix}1& & & & &\\ &\mathrm{det}(g_1)^{-1}\mathrm{det}(g_2)& & & & \\ & & I_2& & &\\  & & &\mathrm{det}(g_1)^{-1}\mathrm{det}(g_2)I_2 & & \\   & & & &1 & \\   & & & & &\mathrm{det}(g_1)^{-1}\mathrm{det}(g_2)\end{smallmatrix}\right)\in  Q_{(2,4)}(\A),\]
which commutes with $\omega_7$. Note
    \[\omega_7 e(\mathrm{diag}(g_1,g_2),\mathrm{det}(g_1)) \alpha \omega_7^{-1} = \left(\begin{smallmatrix}\mathrm{det}(g_2)A'\;^t( g_1 \cdot t_{g_1,g_2} )^{-1}A'^{-1}  & & & \star\\ &g_2& & \\ & &\mathrm{det}(g_2)I_2'\;^tg_2^{-1}I_2'& \\ & & &A g_1 \cdot t_{g_1,g_2} A^{-1}\end{smallmatrix}\right)\in Q_{(2,4)}(\A),\]
    where $A= \left(\begin{smallmatrix}1& \\ &-1\end{smallmatrix}\right)$, $ A':= \left(\begin{smallmatrix} & -1\\1& \end{smallmatrix}\right)$, and  $t_{g_1,g_2} := \left(\begin{smallmatrix}1 & \\ &\mathrm{det}(g_1)^{-1}\mathrm{det}(g_2)\end{smallmatrix}\right)$. Thus, multiplying and dividing by $\alpha$ to the right of $\omega_7e(\mathrm{diag}(g_1,g_2),\mathrm{det}(g_1))\omega_7^{-1}$, by Proposition \ref{Proposition:Pullback:Orthogonal} (1), we have that the integral \eqref{Second:Eq:Prop:Final:Unfolding} equals
        \begin{align*}
     \int_{ \mathbf{Y}(\A)^\circ } \int_{Z_{\GL_4}(\A)\backslash [\GL_2\times \GL_2]} \!\!\!\!\varphi({\rm diag}(g_1,g_2)r) 
       q_{(\overline{\sigma} \boxtimes \sigma)(\tilde{g}_2^{-1}) \cdot \overline{\Psi}'\boxtimes\Psi''}(\Phi_{s})(\omega_7 \alpha^{-1}e(r,a))dg_1dg_2drda,
       \end{align*}
       where we denoted $\tilde{g}_2 := \left(\begin{smallmatrix}g_2& \\ &\mathrm{det}(g_2)I_2'\;^tg_2^{-1}I_2'\end{smallmatrix}\right)$. Since $\alpha^{-1} = e\left(m_{g_1,g_2}, \mathrm{det}(g_1)\mathrm{det}(g_2)^{-1}\right)$, where $$m_{g_1,g_2} := \left(\begin{smallmatrix}1 & & & \\ &\mathrm{det}(g_1)\mathrm{det}(g_2)^{-1}& & \\ & &1 & \\ & & &1\end{smallmatrix}\right)\in \GL_4(\A),$$  we do the change of variables \[\mathrm{det}(g_1)\mathrm{det}(g_2)^{-1}a\mapsto a,\;\;\;\;m_{g_1,g_2}r\mapsto r,\]
       to obtain
       \begin{align*}
       &\int_{ \mathbf{Y}(\A)^\circ } \int_{Z_{\GL_4}(\A)\backslash [\GL_2\times \GL_2]} \!\!\!\!\varphi({\rm diag}(g_1 ,g_2) m_{g_1,g_2}^{-1} r) 
       q_{(\overline{\sigma} \boxtimes \sigma)(\tilde{g}_2^{-1}) \cdot \overline{\Psi}'\boxtimes\Psi''}(\Phi_{s})(\omega_7 e(r,a))dg_1dg_2drda \\ 
        &=\int_{ \mathbf{Y}(\A)^\circ } \int_{Z_{\GL_4}(\A)\backslash [\GL_2\times \GL_2]} \!\!\!\!\varphi({\rm diag}(g_1 t_{g_1,g_2},g_2) r) 
       q_{(\overline{\sigma} \boxtimes \sigma)(\tilde{g}_2^{-1}) \cdot \overline{\Psi}'\boxtimes\Psi''}(\Phi_{s})(\omega_7 e(r,a))dg_1dg_2drda .
        \end{align*}
      Lastly, we do the change of variables $g_1t_{g_1,g_2}\mapsto g_1$, so that the above integral equals 
        \begin{align*}
        \int_{ \mathbf{Y}(\A)^\circ } \int_{Z_{\GL_4}(\A)\backslash [\GL_2\times \GL_2]} \!\!\!\!\varphi({\rm diag}(g_1,g_2) r) 
       q_{(\overline{\sigma} \boxtimes \sigma)(\tilde{g}_2^{-1}) \cdot \overline{\Psi}'\boxtimes\Psi''}(\Phi_{s})(\omega_7 e(r,a))dg_1dg_2drda.  
             \end{align*}
      The result now follows from the fact that, via the exceptional isomorphism, $ \tilde{g}_2$ corresponds to $(I_2, g_2)$ (see Remark \ref{Remark:Siegel:Parabolic:Exceptional}). Then 
    \begin{align*}q_{(\overline{\sigma} \boxtimes \sigma)(\tilde{g}_2^{-1}) \cdot \overline{\Psi}'\boxtimes\Psi''}(\Phi_{s})(\omega_7e(r,a)) = \left( \Lambda_{r,a,s}'\boxtimes\Lambda''_{r,a,s}\right)(\tilde{g_2}) = \Lambda_{r,a,s}'(I_2)\Lambda''_{r,a,s}(g_2),
    \end{align*}
    obtaining the result. 
\end{proof}

\end{document}
